\section{The site-selection problem}
\label{sec:problem}

\subsection{Data requirements for site-selection}
\label{subsec:problem-setup}

We assume the analyst holds data from the $S$ existing sites, which we call
the source $\pool$, and faces a candidate set $\candset=\{1,\ldots,K\}$ of
prospective sites. For each candidate $k$ we know the context profile
$\cand{k}\in\R^p$, which has same basis as the source profiles. This context is
site-level information. For example, in the \citet{meager2019understanding} dataset we use interest rates, program rules, administrative
statistics. Often these data are available from public sources before any experiment is conducted.
We also need some information about the candidate's unit-level covariate
distribution. This can often come from a census, a survey, a prior, or the assumption that units
there resemble units at a reference site. We do not have experimental results from the candidate sites, such as
treatment assignments or outcome data, because no experiment has been
conducted there.

We also state the target we care about as a set $\targset$ of context
profiles. This can be a single profile, a list of benchmark profiles,
or a box around a forecast of the target environment. A set of profiles
is enough, because the width at a target does not depend on its
covariate law (\cref{cor:width-offset}). The design
problem is
\begin{equation}
  k^\star \in \arg\min_{k\in\candset}\;
  \sup_{\ct\in\targset}\;
  \width\bigl(\ct;\ \pool\cup\{k\}\bigr),
  \label{eq:minimax}
\end{equation}
the minimax width of the identified set over the targets, evaluated as
if candidate $k$'s experiment had been run. When candidates carry
costs, as is often the case in practice, the investigator can add a cost term to the objective or screen
the pool by budget first. Everything below goes through in either case because a cost
enters separably in $k$. Choosing a
site is a discrete choice that is only made once, and it should not depend on which single target profile turns out to be the relevant one. Hence, the minimax design is a natural encoding of this problem.

The main obstacle is that $\width(\ct;\pool\cup\{k\})$ depends on data that
we do not have. The following two sections show how much we can determine via the candidate's context alone.

\subsection{What adding a site changes}
\label{subsec:augmentation}

Suppose the candidate experiment is run, with $n_k$ observations and a
randomized treatment within the site, and its data are pooled with the
existing sites. Three inputs of the bound programs change.

First, the span. The centered profile span grows from $\Vspan$ to
\begin{equation}
  \Vspan'=\Vspan+\operatorname{span}\{\cand{k}-c_1\},
  \label{eq:span-update}
\end{equation}
so its dimension rises by one when $\cand{k}-c_1\notin\Vspan$ and is
unchanged otherwise. Notice that the span update is a function of the profiles
alone. That is, we know it at ``site choice'' time.

Second, the moments. The pooled law now mixes in the candidate site, so
$\CovC$, $\dmom{a}$, $\muC$, and hence $\bT$ and $\ellvec$ all move.
Under \cref{ass:model} the new normal equations remain consistent with
the same $\beta_a$. In \cref{sec:selection} we show that their entire
effect on the feasible set is captured by one scalar per arm.

Third, the equality constraint rank. Whether the extra span direction becomes an extra rank
of $\CovC$ depends on covariate overlap. The matrix $\CovC$ averages the
conditional covariance of $C$ given $X$, and a site contributes context
variation only at covariate values where units could have come from more
than one site. A candidate whose covariate support is disjoint from the pool can
contribute profile diversity but no conditional variation. In that case
site membership is a deterministic function of $X$ on the candidate's
support, the residual $\Ctil$ vanishes there, and the added direction is
not identified. We therefore maintain the eigenvalue condition of
\cref{sec:setting} for the augmented pool: the covariance of
$\Bspan'^\top\Ctil$ has a positive smallest eigenvalue, where $\Bspan'$
is a basis of $\Vspan'$. In words, the candidate must overlap the pool
in unit-level covariates, in the same way the target must.

Throughout, the target objects $(\tP,\ct)$ are unchanged by the addition. Only the source pool grows.
